\documentclass[leqno]{article}
% Shared mark scheme preamble for resources.danieljsmith.org: the pack
% preamble, unchanged, then the scheme's own layer. A scheme's .tex is
% rendered from its YAML (scripts/README.md, Mark schemes), never edited.
\input{\CurrentFilePath/pack_preamble.tex}

% The page: equal side margins, so the table runs the full width of the
% header rule. The footers' right edge follows the text's.
\geometry{left=0.8in, right=0.8in}
\makeatletter
\renewcommand{\djs@footR}[1]{%
  \put(\dimexpr1in+\hoffset+\oddsidemargin+\textwidth\relax,
    -\dimexpr\paperheight-0.32in\relax){\makebox[0pt][r]{#1}}}
\makeatother

\usepackage[export]{adjustbox}
\usepackage{tabularray}
\definecolor{msGrid}{RGB}{176, 172, 164}
\definecolor{msTint}{RGB}{245, 243, 239}

% The header on every page, the front page included, and the booklet label
% under the contents-page title. \djsQuestionPageStyle stays empty.
\newcommand{\djsSchemeHeader}[1]{%
  \djsPackHeader[markscheme]{\textit{Mark scheme} - #1}%
  \renewcommand{\djsPackTitle}{#1}%
  \renewcommand{\djsBookletLabel}{Mark scheme}%
}

% The front page: the contents heading, spaced as tocloft spaces the packs'
% "Questions", then the contents table. A question's name and page number
% both link to its first page.
\newcommand{\msContentsHeading}{%
  \vskip\cftbeforetoctitleskip
  {\noindent\cfttoctitlefont\contentsname\par}%
  \vskip\cftaftertoctitleskip}
\NewTblrEnviron{mscontents}
\SetTblrInner[mscontents]{
  colspec = {Q[wd=26mm, l] X[l] Q[wd=16mm, c] Q[wd=14mm, r]},
  width = \linewidth,
  rowsep = 4pt, colsep = 6pt,
  hlines = {0.4pt, msGrid}, vlines = {0.4pt, msGrid},
  hline{2} = {2.2pt, scarlet},
  row{1} = {font=\small\bfseries},
  row{Z} = {font=\bfseries},
  column{2} = {font=\small},
}
\newcommand{\msContentsQuestion}[1]{\hyperref[q:#1]{Question #1}}
\newcommand{\msContentsPage}[1]{\hyperref[q:#1]{\pageref*{q:#1}}}
\newcommand{\msSep}{\hspace{0.9em}}

% The conventions pages' layout: headings in the documents' own bold, tables
% ruled and padded as the scheme's, roomy lists.
\newcommand{\msHead}[1]{%
  \par\vspace{18pt}%
  {\noindent\bfseries #1\par}%
  \vspace{6pt}}
\NewTblrEnviron{mstblr}
\SetTblrInner[mstblr]{rowsep = 6pt, colsep = 6pt,
  hlines = {0.4pt, msGrid}, vlines = {0.4pt, msGrid}}
\newenvironment{msList}
  {\begin{itemize}[leftmargin=1.2em, itemsep=5pt, topsep=3pt, parsep=0pt]}
  {\end{itemize}}

% A question's first page: the link target and the PDF bookmark, as
% \questionitem gives the packs.
\newcommand{\msQnum}{}
\newcommand{\msQuestion}[1]{%
  \clearpage
  \renewcommand{\msQnum}{#1}%
  \phantomsection\label{q:#1}%
  \addcontentsline{toc}{subsection}{\protect\djsTocEntry{#1}}}

% The scheme's table: Question 16mm, Scheme the rest, Marks 16mm, Notes
% 44mm; every cell ruled in msGrid; the scarlet rule under the head row. A
% long table has no caption or "continued" foot; a continuation page opens
% with a muted line above the repeated head row. The rules across a part's
% rows, the alternative's tint and the total row are the rendered table's.
\NewTblrTheme{ms}{%
  \DefTblrTemplate{firsthead}{default}{}%
  \DefTblrTemplate{middlehead,lasthead}{default}{%
    {\small\itshape\color{djsMuted}Question \msQnum\ continued\par}%
    \vspace{6pt}}%
  \DefTblrTemplate{firstfoot,middlefoot,lastfoot}{default}{}%
}
\NewTblrEnviron{msscheme}
\SetTblrOuter[msscheme]{long, theme=ms, entry=none, label=none}
\SetTblrInner[msscheme]{
  colspec = {Q[wd=16mm, l] X[l] Q[wd=16mm, c] Q[wd=44mm, l]},
  width = \linewidth,
  rowsep = 6pt, colsep = 6pt,
  rowhead = 1,
  vlines = {0.4pt, msGrid},
  hline{1} = {0.4pt, msGrid}, hline{2} = {2.2pt, scarlet},
  hline{Z} = {0.4pt, msGrid},
  column{1} = {font=\bfseries}, column{3} = {font=\bfseries},
  column{4} = {font=\small},
  row{1} = {rowsep=4pt, font=\small\bfseries},
}

% Two marks on a row: each note starts with its code.
\newcommand{\msNote}[2]{\textbf{#1}\enspace #2}
% A full-row figure: the worked solution's own picture, scaled uniformly and
% only down, to the table less the Question column. The picture is drawn
% with the line width it has in the solutions document, at its list depth
% there (the optional argument, 1 by default: a part's working), so one
% sized by \linewidth keeps its size; tabularray sets the cell's \linewidth
% to the cell's measured width, which would shrink it at every pass. The
% widths are measured once, in a solution box at the packs' text width
% (pack_preamble.tex's margins, 0.8in and 1in).
\newlength{\msSpanW}
\makeatletter
\newcommand{\ms@keepLW}[1]{%
  \expandafter\xdef\csname ms@solnLW@#1\endcsname{\the\linewidth}}
\AtBeginDocument{%
  \setlength{\msSpanW}{\dimexpr\textwidth-16mm-28pt\relax}%
  \setbox\z@\vbox{%
    \begin{minipage}{\dimexpr\paperwidth-1.8in\relax}
      \begin{djsSolution}
        \ms@keepLW{0}%
        \begin{enumerate}\item \ms@keepLW{1}%
        \begin{enumerate}\item \ms@keepLW{2}%
        \begin{enumerate}\item \ms@keepLW{3}%
        \end{enumerate}\end{enumerate}\end{enumerate}
      \end{djsSolution}
    \end{minipage}}}
\newcommand{\msFigure}[2][1]{%
  \adjustbox{max width=\msSpanW, max totalheight=0.6\textheight, valign=t}{%
    \hsize=\csname ms@solnLW@#1\endcsname
    \linewidth=\hsize
    #2}}
\makeatother
% A figure mark's feature, in the Scheme column.
\newcommand{\msFeature}[1]{\textit{#1}}
% An alternative's title row.
\newcommand{\msAlternative}[1]{\textit{Alternative:} #1}
% Drafts only: a queried mark's code, the query's why under its note, and a
% flag in its part's place.
\newcommand{\msQueried}[2]{#1~\(\to\) #2?}
\newcommand{\msQuery}[1]{{\itshape\color{djsMuted}Query: #1}}
\newcommand{\msFlag}[2]{\textbf{#1:}\enspace #2}


\djsSchemeHeader{De Moivre's Theorem}
\djsSitePath{Core Pure}{Complex Numbers}

\djsLastUpdated{26 October 2026}
\begin{document}
\djsSetQuestionCount{18}
\djsAddMarks{1}{5}
\djsAddMarks{2}{7}
\djsAddMarks{3}{7}
\djsAddMarks{4}{9}
\djsAddMarks{5}{11}
\djsAddMarks{6}{9}
\djsAddMarks{7}{8}
\djsAddMarks{8}{9}
\djsAddMarks{9}{10}
\djsAddMarks{10}{10}
\djsAddMarks{11}{12}
\djsAddMarks{12}{12}
\djsAddMarks{13}{9}
\djsAddMarks{14}{11}
\djsAddMarks{15}{9}
\djsAddMarks{16}{19}
\djsAddMarks{17}{13}
\djsAddMarks{18}{14}
\vspace*{6pt}
\djsContentsTitle
\vspace{30pt}
\msContentsHeading
\begin{mscontents}{}
Question & Marks by part & Marks & Page \\
\msContentsQuestion{1} & & 5 & \msContentsPage{1} \\
\msContentsQuestion{2} & (a)\,3\msSep (b)\,4 & 7 & \msContentsPage{2} \\
\msContentsQuestion{3} & & 7 & \msContentsPage{3} \\
\msContentsQuestion{4} & (a)\,5\msSep (b)\,4 & 9 & \msContentsPage{4} \\
\msContentsQuestion{5} & (a)\,5\msSep (b)\,6 & 11 & \msContentsPage{5} \\
\msContentsQuestion{6} & (a)\,5\msSep (b)\,4 & 9 & \msContentsPage{6} \\
\msContentsQuestion{7} & (a)\,3\msSep (b)\,5 & 8 & \msContentsPage{7} \\
\msContentsQuestion{8} & (a)\,3\msSep (b)\,4\msSep (c)\,2 & 9 & \msContentsPage{8} \\
\msContentsQuestion{9} & (a)\,6\msSep (b)\,4 & 10 & \msContentsPage{9} \\
\msContentsQuestion{10} & (a)\,5\msSep (b)\,5 & 10 & \msContentsPage{10} \\
\msContentsQuestion{11} & (a)\,5\msSep (b)\,4\msSep (c)\,3 & 12 & \msContentsPage{11} \\
\msContentsQuestion{12} & (i)\,5\msSep (ii)\,7 & 12 & \msContentsPage{12} \\
\msContentsQuestion{13} & (a)\,5\msSep (b)\,4 & 9 & \msContentsPage{13} \\
\msContentsQuestion{14} & (a)\,5\msSep (b)\,6 & 11 & \msContentsPage{14} \\
\msContentsQuestion{15} & (a)\,5\msSep (b)\,4 & 9 & \msContentsPage{15} \\
\msContentsQuestion{16} & (i)(A)\,2\msSep (i)(B)\,2\msSep (ii)\,6\msSep (iii)\,6\msSep (iv)\,3 & 19 & \msContentsPage{16} \\
\msContentsQuestion{17} & (a)\,5\msSep (b)\,3\msSep (c)\,5 & 13 & \msContentsPage{17} \\
\msContentsQuestion{18} & (a)\,5\msSep (b)\,6\msSep (c)\,3 & 14 & \msContentsPage{18} \\
\SetCell[c=2]{r} Total & & 184 & \\
\end{mscontents}
\newpage
% The conventions pages, shared by every pack's mark scheme: one source of
% truth, \input after the front page, with no strand or pack logic. The keys
% on the first side, the rules on the second. The scheme layer defines
% \msHead, msList and mstblr.
\vspace*{6pt}%
{\noindent\cfttoctitlefont Reading the mark scheme\par}%
\vspace{10pt}%
These rules apply throughout, unless a note says otherwise.

\msHead{Marks}
Every mark is worth one, so \textbf{M1 A1} on a row is two marks. Only marks
coded ft follow through.\par
\vspace{10pt}%
\begin{mstblr}{
  colspec = {Q[wd=15mm, l] X[l]},
  width = \linewidth,
  hline{2} = {2.2pt, scarlet},
  column{1} = {font=\bfseries},
  row{1} = {font=\small\bfseries, rowsep=4pt},
}
Code & Meaning \\
M1 & Method: a correct method, applied to this question. A slip doesn't cost it if the method is clear. A formula must be used, not just quoted. \\
dM1 & Dependent method: as M1, but only if the method mark before it was earned, unless its note names another. \\
A1 & Accuracy: a correct answer or step. It needs the method mark it follows. \\
A1ft & As A1, following through from an earlier answer that the note names. \\
A1~AG & The last mark of a part whose answer is printed in the question: complete working, with no errors. \\
B1 & A correct value, fact or statement. It needs no method mark. \\
B1ft & As B1, following through from an earlier answer that the note names. \\
E1 & A written explanation, interpretation or conclusion. Concluding a proof by induction needs every earlier mark. \\
\end{mstblr}

\msHead{Abbreviations}
ft and AG qualify a code in the Marks column. The others appear in the
notes.\par
\vspace{10pt}%
\begin{mstblr}{
  colspec = {Q[wd=17mm, l] X[l] Q[wd=15mm, l] X[l]},
  width = \linewidth,
  column{1} = {font=\bfseries}, column{3} = {font=\bfseries},
}
ft & follow through & AG & answer given in the question \\
cao & correct answer only, with no follow-through & cso & correct solution only: no errors in the part \\
oe & or an equivalent form & awrt & any answer that rounds to this value \\
soi & seen, or implied by later working & isw & ignore working after a correct answer \\
SC & special case: credit outside the main method & their & the student's own earlier value \\
condone & a slip that doesn't cost the mark & [\;\,] & working the mark doesn't need \\
s.f. & significant figures & d.p. & decimal places \\
\end{mstblr}

\newpage
\msHead{Reading the scheme}
\begin{msList}
\item Each row is a line of working, with the marks it earns beside it.
  Where a row carries two marks, each note starts with its code.
\item \textbf{[\textit{n}]} ends each part with its total.
\item A shaded \textit{Alternative} is a second method for the same total.
  Mark an answer by the method it follows.
\item Any other valid method earns the equivalent marks, unless the question
  names a method or says \textit{hence}.
\item A sketch earns one B1 for each feature named below its figure.
\end{msList}

\msHead{Accuracy and answers}
\begin{msList}
\item Unless a question says otherwise, give answers to 3 significant
  figures, and angles in degrees to 1 decimal place. Angles in radians take
  3 significant figures.
\item A final answer with more figures is accepted if it rounds to the
  scheme's value; one with fewer is not. An answer spoiled by rounding too
  early loses its accuracy mark.
\item Where a question asks for an exact answer, a decimal loses the final
  accuracy mark. Otherwise either is accepted.
\item Wrong or missing units cost nothing unless the question asks for them.
\item A correct answer with no working earns full marks, unless the question
  needs the working: show that, prove, hence, a method it names, or an exact
  answer. A wrong answer with no working earns nothing. Calculator answers
  are treated the same way.
\item A correct answer, once reached, isn't spoiled by later working (isw).
\end{msList}

\msHead{Value of $\boldsymbol{g}$}
\begin{msList}
\item Take $g = 9.8\,\mathrm{m\,s^{-2}}$ unless a question states a value,
  which is then the only one accepted. Answers that use $g = 9.8$ may be
  given to 2 or 3 significant figures.
\item Where a question doesn't state $g$, an answer worked with $g = 10$
  throughout earns the same marks, keeping any value printed in the
  question. The worked solutions list the final answers with $g = 10$.
\end{msList}


\msQuestion{1}
\begin{msscheme}{
  hline{3-6} = {2-Z}{0.4pt, msGrid},
  hline{7} = {0.4pt, msGrid},
  row{7} = {font=\bfseries},
}
Question & Scheme & Marks & Notes \\
1 & {\(z - z^{-1} = 2i\sin\theta,\ z + z^{-1} = 2\cos\theta\)\\[4pt] \(\Rightarrow \sin\theta\cos^4\theta = \frac{1}{32i}(z-z^{-1})(z+z^{-1})^4\)} & B1 & {Both, stated or used.} \\
& {\((z-z^{-1})(z^4+4z^2+6+4z^{-2}+z^{-4})\)} & M1 & {Multiplies out \((z-z^{-1})(z+z^{-1})^4\) to a sum of powers of \(z\), in any order of multiplication. Condone coefficient and sign slips.} \\
& {\(= z^5+3z^3+2z-2z^{-1}-3z^{-3}-z^{-5}\)} & A1 & {Correct expansion, oe; the factor \(\frac{1}{32i}\) need not be shown here.} \\
& {\((z^5-z^{-5})+3(z^3-z^{-3})+2(z-z^{-1})\)\\[4pt] \(\phantom{(} = 2i\sin5\theta+6i\sin3\theta+4i\sin\theta\)} & M1 & {Pairs terms as \(z^n - z^{-n}\) and replaces each by \(2i\sin n\theta\), keeping their coefficients.} \\
& {\(A=\frac18,\ B=\frac{3}{16},\ C=\frac{1}{16}\)} & A1 & {cso; withhold if an \(i\) or the factor \(\frac{1}{32}\) is lost and recovered. Accept the identity written with these values.} \\
\nopagebreak \SetCell[c=2]{r} Total for Question 1 & & 5 & \\
\end{msscheme}

\msQuestion{2}
\begin{msscheme}{
  hline{3-5,7-10} = {2-Z}{0.4pt, msGrid},
  hline{6,11} = {0.4pt, msGrid},
  row{11} = {font=\bfseries},
}
Question & Scheme & Marks & Notes \\
2(a) & {\(z^n = \cos n\theta + i\sin n\theta\)\\[4pt] \(z^{-n} = \cos(-n\theta) + i\sin(-n\theta)\)\\[4pt] \(\phantom{z^{-n}} = \cos n\theta - i\sin n\theta\)} & M1 & {Uses de Moivre's theorem to write both \(z^n\) and \(z^{-n}\) in terms of \(\cos n\theta\) and \(\sin n\theta\); \(z^{-n}\) may come from \(\frac{1}{z^n}\) by realising the denominator.} \\
& {\(z^n + \frac{1}{z^n} = 2\cos n\theta\)} & A1 & {From correct expressions for both powers, with \(\cos(-n\theta)\) and \(\sin(-n\theta)\) correctly simplified.} \\
& {\(z^n - \frac{1}{z^n} = (\cos n\theta + i\sin n\theta) - (\cos n\theta - i\sin n\theta)\)\\[4pt] \(\phantom{z^n - \frac{1}{z^n}} = 2i\sin n\theta\)} & A1~AG & {Correct subtraction of the two expressions leading without error to the printed result.} \\
\nopagebreak & & [3] & \\
2(b) & {\(\left(z+\frac1z\right)^3 = z^3+\frac{1}{z^3} + 3\left(z+\frac1z\right)\)\\[4pt] \(\Rightarrow 8\cos^3\theta = 2\cos3\theta + 6\cos\theta\)\\[4pt] \(\left(z-\frac1z\right)^3 = z^3-\frac{1}{z^3} - 3\left(z-\frac1z\right)\)\\[4pt] \(\Rightarrow -8i\sin^3\theta = 2i\sin3\theta - 6i\sin\theta\)} & M1 & {Expands both cubes binomially, groups into \(z^3 \pm z^{-3}\) and \(z \pm z^{-1}\), and replaces each by (a). Condone sign and coefficient slips.} \\
& {\(\cos3\theta = 4\cos^3\theta - 3\cos\theta\)\\[4pt] \(\sin3\theta = 3\sin\theta - 4\sin^3\theta\)} & A1 & {Both correct, from working with \((2i\sin\theta)^3 = -8i\sin^3\theta\) correct.} \\
& {\(\tan3\theta = \frac{3\sin\theta - 4\sin^3\theta}{4\cos^3\theta - 3\cos\theta}\)\\[4pt] \(\phantom{\tan3\theta} = \frac{3\tan\theta\sec^2\theta - 4\tan^3\theta}{4 - 3\sec^2\theta}\)} & M1 & {Writes \(\tan3\theta\) as their \(\sin3\theta\) over their \(\cos3\theta\), each a cubic in \(\sin\theta\) or \(\cos\theta\), and divides through by a power of \(\cos\theta\) to introduce \(\tan\theta\).} \\
& {\(\phantom{\tan3\theta} = \frac{3\tan\theta(1+\tan^2\theta) - 4\tan^3\theta}{4 - 3(1+\tan^2\theta)}\)\\[4pt] \(\phantom{\tan3\theta} = \frac{3\tan\theta - \tan^3\theta}{1 - 3\tan^2\theta}\)} & A1~AG & {From correct expressions for \(\sin3\theta\) and \(\cos3\theta\), with \(\sec^2\theta = 1+\tan^2\theta\) oe used correctly, leading without error to the printed result.} \\
\nopagebreak & & [4] & \\
\nopagebreak \SetCell[c=2]{r} Total for Question 2 & & 7 & \\
\end{msscheme}

\msQuestion{3}
\begin{msscheme}{
  hline{3-7} = {2-Z}{0.4pt, msGrid},
  hline{8} = {0.4pt, msGrid},
  row{8} = {font=\bfseries},
}
Question & Scheme & Marks & Notes \\
3 & {\(\cos 6\theta + i\sin 6\theta = (\cos\theta + i\sin\theta)^6\)\\[4pt] \(\phantom{\cos 6\theta + i\sin 6\theta} = \cos^6\theta\,(1+it)^6\)} & M1 & {Applies de Moivre's theorem and equates \(\cos 6\theta + i\sin 6\theta\) to \(\cos^6\theta\,(1+it)^6\).} \\
& {\((1+it)^6 = 1 + 6(it) + 15(it)^2 + 20(it)^3\)\\[4pt] \(\phantom{(1+it)^6 =} + 15(it)^4 + 6(it)^5 + (it)^6\)} & M1 & {Expands with seven terms in powers of \(it\) and binomial coefficients; condone one coefficient slip.} \\
& {\(\phantom{(1+it)^6} = (1 - 15t^2 + 15t^4 - t^6)\)\\[4pt] \(\phantom{(1+it)^6 =} + i(6t - 20t^3 + 6t^5)\)} & A1 & {Correct, with powers of \(i\) simplified; need not be grouped into real and imaginary parts.} \\
& {\(\cos 6\theta = \cos^6\theta\,(1 - 15t^2 + 15t^4 - t^6)\)\\[4pt] \(\sin 6\theta = \cos^6\theta\,(6t - 20t^3 + 6t^5)\)} & M1 A1 & {\msNote{M1}{Equates real parts and imaginary parts of their two expressions for \((\cos\theta + i\sin\theta)^6\).}\\ \msNote{A1}{Both correct, including the factor \(\cos^6\theta\).}} \\
& {\(\tan 6\theta = \frac{\sin 6\theta}{\cos 6\theta} = \frac{\cos^6\theta\,(6t - 20t^3 + 6t^5)}{\cos^6\theta\,(1 - 15t^2 + 15t^4 - t^6)}\)} & M1 & {Divides their \(\sin 6\theta\) by their \(\cos 6\theta\), or the imaginary part by the real part, and cancels \(\cos^6\theta\).} \\
& {\(\tan 6\theta = \frac{6\tan\theta - 20\tan^3\theta + 6\tan^5\theta}{1 - 15\tan^2\theta + 15\tan^4\theta - \tan^6\theta}\)} & A1~AG & {cso: printed result from correct working with no errors in powers of \(i\) or signs. Condone the final line left in \(t\).} \\
\nopagebreak \SetCell[c=2]{r} Total for Question 3 & & 7 & \\
\end{msscheme}

\msQuestion{4}
\begin{msscheme}{
  hline{3-7,9-12} = {2-Z}{0.4pt, msGrid},
  hline{8,13} = {0.4pt, msGrid},
  row{13} = {font=\bfseries},
}
Question & Scheme & Marks & Notes \\
4(a) & {\((\cos x + i\sin x)^6 = \cos 6x + i\sin 6x\)\\[4pt] \(\phantom{(\cos x + i\sin x)^6} = \cos^6 x + 6i\cos^5 x\sin x\)\\[4pt] \(\phantom{(\cos x + i\sin x)^6 =} + 15i^2\cos^4 x\sin^2 x + \ldots\)} & M1 & {Applies de Moivre to \((\cos x + i\sin x)^6\) and expands binomially, at least the terms giving the imaginary part. Condone slips in coefficients or powers of \(i\).} \\
& {\(\sin 6x = 6\cos^5 x\sin x - 20\cos^3 x\sin^3 x\)\\[4pt] \(\phantom{\sin 6x =} + 6\cos x\sin^5 x\)} & A1 & {Equates imaginary parts correctly, oe.} \\
& {\(\sin 6x = 2\sin x\cos x\big(3\cos^4 x\)\\[4pt] \(\phantom{\sin 6x =} - 10\cos^2 x(1-\cos^2 x) + 3(1-\cos^2 x)^2\big)\)} & M1 & {Uses \(\sin^2 x + \cos^2 x = 1\) to write the bracket in powers of \(\cos x\) alone, by any route.} \\
& {\(\sin 6x = \sin 2x\,(\ldots\cos^4 x + \ldots\cos^2 x + \ldots)\)} & dM1 & {Expands and collects to a quartic in \(\cos x\), with \(2\sin x\cos x\) replaced by \(\sin 2x\). Condone arithmetic slips.} \\
& {\(\sin 6x = \sin 2x\,(16\cos^4 x - 16\cos^2 x + 3)\)\\[4pt] \(a = 16,\ b = -16,\ c = 3\)} & A1 & {cso. Correct identity or \(a\), \(b\), \(c\) stated.} \\
\nopagebreak & & [5] & \\
4(b) & {\(\sin 2\theta(16\cos^4\theta - 16\cos^2\theta + 3) = \sin 2\theta\cos 2\theta\)\\[4pt] \(\Rightarrow 16\cos^4\theta - 16\cos^2\theta + 3 = 2\cos^2\theta - 1\)} & M1 & {Uses their identity of the form \(\sin 2\theta(a\cos^4\theta + b\cos^2\theta + c)\), cancels \(\sin 2\theta\) and writes \(\cos 2\theta\) as \(2\cos^2\theta - 1\). Condone sign slips.} \\
& {\(16\cos^4\theta - 18\cos^2\theta + 4 = 0\)} & A1 & {oe, such as \(8\cos^4\theta - 9\cos^2\theta + 2 = 0\).} \\
& {\(\cos^2\theta = \frac{9 \pm \sqrt{17}}{16} \Rightarrow \theta = \ldots\)} & M1 & {Solves their three-term quadratic in \(\cos^2\theta\) by any means, calculator allowed, and inverts to find at least one \(\theta\).} \\
& {\(\theta = 0.438,\ 0.986\)} & A1 & {Both, awrt 0.438 and 0.986, and no others in \(0 < \theta < \frac{\pi}{2}\).} \\
\nopagebreak & & [4] & \\
\nopagebreak \SetCell[c=2]{r} Total for Question 4 & & 9 & \\
\end{msscheme}

\msQuestion{5}
\begin{msscheme}{
  hline{3-6,8-12} = {2-Z}{0.4pt, msGrid},
  hline{7,13} = {0.4pt, msGrid},
  row{13} = {font=\bfseries},
}
Question & Scheme & Marks & Notes \\
5(a) & {\((z - z^{-1})^6 = (2i\sin\theta)^6 = -64\sin^6\theta\)} & B1 & \\
& {\((z - z^{-1})^6\)\\[4pt] \(\phantom{A} = z^6 - 6z^4 + 15z^2 - 20 + 15z^{-2} - 6z^{-4} + z^{-6}\)} & M1 A1 & {\msNote{M1}{Expands using binomial coefficients 1, 6, 15, 20 with powers of \(z\) and \(z^{-1}\). Condone sign and power slips; need not be simplified.}\\ \msNote{A1}{All seven terms correct; powers need not be simplified.}} \\
& {\((z^6 + z^{-6}) - 6(z^4 + z^{-4}) + 15(z^2 + z^{-2}) - 20\)\\[4pt] \(\phantom{A} = 2\cos 6\theta - 12\cos 4\theta + 30\cos 2\theta - 20\)} & M1 & {Pairs terms and uses \(z^n + z^{-n} = 2\cos n\theta\) for at least two pairs. Condone coefficient slips.} \\
& {\(-64\sin^6\theta = 2\cos 6\theta - 12\cos 4\theta + 30\cos 2\theta - 20\)\\[4pt] \(\Rightarrow \sin^6\theta = \frac{1}{32}(10 - 15\cos 2\theta + 6\cos 4\theta - \cos 6\theta)\)} & A1~AG & {Equates the two expressions and reaches the printed identity without error; combined rearrangement accepted.} \\
\nopagebreak & & [5] & \\
5(b) & {\(x = \tan\theta,\ \mathrm{d}x = \sec^2\theta\,\mathrm{d}\theta\)\\[4pt] \(\Rightarrow \int \frac{\tan^6\theta}{(\sec^2\theta)^4}\sec^2\theta\,\mathrm{d}\theta\)} & M1 & {Substitutes fully: \(\mathrm{d}x\) replaced by \(\sec^2\theta\,\mathrm{d}\theta\) and \(1 + x^2\) by \(\sec^2\theta\) or \(1 + \tan^2\theta\). Condone slips in powers.} \\
& {\(\int_0^{\pi/4} \sin^6\theta\,\mathrm{d}\theta\)} & A1 & \\
& {\(\frac{1}{32}\int_0^{\pi/4} (10 - 15\cos 2\theta + 6\cos 4\theta - \cos 6\theta)\,\mathrm{d}\theta\)\\[4pt] \(\phantom{A} = \frac{1}{32}\left[10\theta - \frac{15}{2}\sin 2\theta + \frac{3}{2}\sin 4\theta - \frac{1}{6}\sin 6\theta\right]_0^{\pi/4}\)} & M1 A1 & {\msNote{M1}{Replaces \(\sin^6\theta\) by the result of (a) and integrates to the form \(a\theta + \sum b_k \sin k\theta\). Condone sign and coefficient slips.}\\ \msNote{A1}{Correct integral; the factor \(\frac{1}{32}\) may be applied later.}} \\
& {\(\frac{1}{32}\left(\frac{10\pi}{4} - \frac{15}{2} + 0 + \frac{1}{6}\right)\)} & dM1 & {Substitutes \(\frac{\pi}{4}\) and 0 into their integral and subtracts; the zero from the lower limit need not be shown.} \\
& {\(= \frac{15\pi - 44}{192}\)} & A1 & {Any exact equivalent, such as \(\frac{5\pi}{64} - \frac{11}{48}\).} \\
\nopagebreak & & [6] & \\
\nopagebreak \SetCell[c=2]{r} Total for Question 5 & & 11 & \\
\end{msscheme}

\msQuestion{6}
\begin{msscheme}{
  hline{3-7,9-12} = {2-Z}{0.4pt, msGrid},
  hline{8,13} = {0.4pt, msGrid},
  row{13} = {font=\bfseries},
}
Question & Scheme & Marks & Notes \\
6(a) & {\(\cos5\theta + i\sin5\theta = (\cos\theta + i\sin\theta)^5\)\\[4pt] \(\phantom{\cos5\theta} = \cos^5\theta + 5i\cos^4\theta\sin\theta - 10\cos^3\theta\sin^2\theta + \ldots\)} & M1 & {Uses de Moivre with a binomial expansion of \((\cos\theta + i\sin\theta)^5\); at least the real terms, with binomial coefficients. Condone slips in powers of \(i\).} \\
& {\(\cos5\theta = \cos^5\theta - 10\cos^3\theta\sin^2\theta + 5\cos\theta\sin^4\theta\)} & A1 & {Equates real parts correctly.} \\
& {\(\phantom{\cos5\theta} = \cos^5\theta - 10\cos^3\theta(1-\cos^2\theta)\)\\[4pt] \(\phantom{\cos5\theta =} + 5\cos\theta(1-\cos^2\theta)^2\)} & M1 & {Replaces \(\sin^2\theta\) by \(1-\cos^2\theta\) in both terms containing \(\sin\theta\).} \\
& {\(\phantom{\cos5\theta} = \cos^5\theta - 10\cos^3\theta + 10\cos^5\theta\)\\[4pt] \(\phantom{\cos5\theta =} + 5\cos\theta - 10\cos^3\theta + 5\cos^5\theta\)} & dM1 & {Expands the brackets and collects to a form \(a\cos^5\theta + b\cos^3\theta + c\cos\theta\). Condone arithmetic slips.} \\
& {\(\cos5\theta = 16\cos^5\theta - 20\cos^3\theta + 5\cos\theta\)\\[4pt] \(2\cos5\theta = 32\cos^5\theta - 40\cos^3\theta + 10\cos\theta\)} & A1~AG & {Correct real part, substituted and expanded without error to the printed identity, including the doubling; combined steps accepted.} \\
\nopagebreak & & [5] & \\
6(b) & {\(x = 2\cos\theta \Rightarrow x^5 - 5x^3 + 5x\)\\[4pt] \(\phantom{x} = 32\cos^5\theta - 40\cos^3\theta + 10\cos\theta = 2\cos5\theta\)} & M1 & {Substitutes \(x = 2\cos\theta\) and uses (a) to write \(x^5 - 5x^3 + 5x\) as \(2\cos5\theta\).} \\
& {\(2\cos5\theta + 1 = 0 \Rightarrow \cos5\theta = -\frac12\)} & A1 & \\
& {\(5\theta = \frac{2\pi}{3}, \frac{4\pi}{3}, \frac{8\pi}{3}, \frac{10\pi}{3}, \frac{14\pi}{3} \Rightarrow \theta = \ldots\)} & M1 & {Finds at least two values of \(5\theta\) for their \(\cos5\theta = k\), \(|k| \le 1\), and divides by 5.} \\
& {\(x = 2\cos\frac{2\pi}{15}, 2\cos\frac{4\pi}{15}, 2\cos\frac{8\pi}{15}, 2\cos\frac{2\pi}{3}, 2\cos\frac{14\pi}{15}\)} & A1 & {All five, each as \(2\cos q\pi\) oe (e.g. \(2\cos\frac{16\pi}{15}\)), no repeated root. Decimals, \(-1\), or \(\theta\) values alone score A0.} \\
\nopagebreak & & [4] & \\
\nopagebreak \SetCell[c=2]{r} Total for Question 6 & & 9 & \\
\end{msscheme}

\msQuestion{7}
\begin{msscheme}{
  hline{3-5,7-11} = {2-Z}{0.4pt, msGrid},
  hline{6,12} = {0.4pt, msGrid},
  row{12} = {font=\bfseries},
}
Question & Scheme & Marks & Notes \\
7(a) & {\(z^n = \cos n\theta + i\sin n\theta\)\\[4pt] \(z^{-n} = \cos(-n\theta) + i\sin(-n\theta)\)\\[4pt] \(\phantom{z^{-n}} = \cos n\theta - i\sin n\theta\)} & M1 & {Applies de Moivre to both \(z^n\) and \(z^{-n}\), or finds \(\frac{1}{z^n}\) from \(z^n\) using the conjugate. Simplification of \(\cos(-n\theta)\), \(\sin(-n\theta)\) not needed.} \\
& {\(z^n + \frac{1}{z^n} = 2\cos n\theta\)} & A1 & {From a correctly simplified \(\frac{1}{z^n}\).} \\
& {\(z^n - \frac{1}{z^n} = 2i\sin n\theta\)} & A1~AG & {Correct subtraction leading without error to the printed result.} \\
\nopagebreak & & [3] & \\
7(b) & {\((2\cos\theta)^3(2i\sin\theta)^2 = -32\sin^2\theta\cos^3\theta\)} & B1 & \\
& {\(\left(z^3+3z+\frac{3}{z}+\frac{1}{z^3}\right)\left(z^2-2+\frac{1}{z^2}\right) = \ldots\)} & M1 & {Multiplies out the product, by binomial expansions or otherwise, to a sum of powers of \(z\). Condone coefficient slips.} \\
& {\(z^5 + z^3 - 2z - \frac{2}{z} + \frac{1}{z^3} + \frac{1}{z^5}\)} & A1 & \\
& {\(\left(z^5+\frac{1}{z^5}\right)+\left(z^3+\frac{1}{z^3}\right)-2\left(z+\frac{1}{z}\right)\)\\[4pt] \(\phantom{A} = 2\cos 5\theta + 2\cos 3\theta - 4\cos\theta\)} & M1 & {Pairs \(z^k\) with \(z^{-k}\) in their expansion and replaces each pair by \(2\cos k\theta\).} \\
& {\(A = -\frac{1}{16},\ B = -\frac{1}{16},\ C = \frac{1}{8}\)} & A1 & {cso, from equating with \(-32\sin^2\theta\cos^3\theta\). Accept the identity written with these values.} \\
\nopagebreak & & [5] & \\
\nopagebreak \SetCell[c=2]{r} Total for Question 7 & & 8 & \\
\end{msscheme}

\msQuestion{8}
\begin{msscheme}{
  hline{3-5,7-10,12-13} = {2-Z}{0.4pt, msGrid},
  hline{6,11,14} = {0.4pt, msGrid},
  row{14} = {font=\bfseries},
}
Question & Scheme & Marks & Notes \\
8(a) & {\((1+i\tan\theta)^4 = \sec^4\theta(\cos\theta+i\sin\theta)^4\)} & B1 & {The given form raised to the fourth power, \(\sec^4\theta(\cos\theta+i\sin\theta)^4\) oe. Not for the given form alone.} \\
& {\(\phantom{(1+i\tan\theta)^4} = \sec^4\theta(\cos4\theta+i\sin4\theta)\)} & M1 & {Applies de Moivre's theorem to \((\cos\theta+i\sin\theta)^4\). Condone a slip in the power of \(\sec\theta\).} \\
& {\(\phantom{(1+i\tan\theta)^4} = \frac{\cos4\theta+i\sin4\theta}{\cos^4\theta}\)} & A1~AG & {Printed result reached without error from the fourth power of the given form and de Moivre's theorem.} \\
\nopagebreak & & [3] & \\
8(b) & {\((1+i\tan\theta)^4 = 1+4i\tan\theta+6i^2\tan^2\theta\)\\[4pt] \(\phantom{(1+i\tan\theta)^4 =} + 4i^3\tan^3\theta + i^4\tan^4\theta\)} & M1 & {Expands to five terms with coefficients 1, 4, 6, 4, 1 and powers of \(i\tan\theta\). Powers of \(i\) need not be simplified.} \\
& {\(= (1-6\tan^2\theta+\tan^4\theta)\)\\[4pt] \(\phantom{=} + i\,4\tan\theta(1-\tan^2\theta)\)} & A1 & {Correct real and imaginary parts, in any arrangement.} \\
& {Equates parts with (a):\\[4pt] \(\frac{\cos4\theta}{\cos^4\theta} = 1-6\tan^2\theta+\tan^4\theta\)\\[4pt] \(\cos^4\theta = \frac{1}{(1+\tan^2\theta)^2}\)} & M1 & {Equates real or imaginary parts with (a) and replaces \(\cos^4\theta\) using \(\sec^2\theta=1+\tan^2\theta\). Condone slips in their expansion.} \\
& {\(\cos4\theta = \frac{1-6\tan^2\theta+\tan^4\theta}{(1+\tan^2\theta)^2}\)\\[4pt] \(\sin4\theta = \frac{4\tan\theta(1-\tan^2\theta)}{(1+\tan^2\theta)^2}\)} & A1~AG & {Both printed results, reached without error from correct real and imaginary parts.} \\
\nopagebreak & & [4] & \\
8(c) & {\(\tan4\theta = \frac{\sin4\theta}{\cos4\theta}\)\\[4pt] \(\phantom{\tan4\theta} = \frac{4\tan\theta(1-\tan^2\theta)}{(1+\tan^2\theta)^2}\)\\[4pt] \(\phantom{\tan4\theta =} \div \frac{1-6\tan^2\theta+\tan^4\theta}{(1+\tan^2\theta)^2}\)} & M1 & {Divides their \(\sin4\theta\) by their \(\cos4\theta\) from (b), or the imaginary part of the expansion by the real part.} \\
& {\(\tan4\theta = \frac{4\tan\theta(1-\tan^2\theta)}{1-6\tan^2\theta+\tan^4\theta}\)} & A1~AG & {Printed result from the correct quotient with no errors.} \\
\nopagebreak & & [2] & \\
\nopagebreak \SetCell[c=2]{r} Total for Question 8 & & 9 & \\
\end{msscheme}

\msQuestion{9}
\begin{msscheme}{
  hline{3-8,10-13} = {2-Z}{0.4pt, msGrid},
  hline{9,14} = {0.4pt, msGrid},
  row{14} = {font=\bfseries},
}
Question & Scheme & Marks & Notes \\
9(a) & {\(\cos 6\theta + i\sin 6\theta = (c + is)^6\)\\[4pt] \(\phantom{\cos} = c^6 + 6c^5(is) + 15c^4(is)^2 + \ldots + (is)^6\)} & M1 & {Equates \(\cos 6\theta + i\sin 6\theta\) to a binomial expansion of \((\cos\theta + i\sin\theta)^6\) with at least the real terms. Condone slips in coefficients or powers of \(i\).} \\
& {\(\cos 6\theta = c^6 - 15c^4s^2 + 15c^2s^4 - s^6\)} & A1 & \\
& {\(\cos 6\theta = c^6 - 15c^4(1-c^2) + 15c^2(1-c^2)^2\)\\[4pt] \(\phantom{\cos 6\theta =} - (1-c^2)^3 = \ldots\)} & M1 & {Replaces every \(\sin^2\theta\) by \(1 - \cos^2\theta\) and expands to a polynomial in \(\cos\theta\). Condone expansion slips.} \\
& {\(\cos 6\theta = 32c^6 - 48c^4 + 18c^2 - 1\)} & A1 & \\
& {\(\sec 6\theta = \dfrac{1}{\their{32c^6 - 48c^4 + 18c^2 - 1}} \times \dfrac{\sec^6\theta}{\sec^6\theta}\)} & M1 & {Takes the reciprocal of their polynomial in \(\cos\theta\) and multiplies numerator and denominator by \(\sec^6\theta\), or writes each \(\cos^n\theta\) as \(\sec^{-n}\theta\).} \\
& {\(\sec 6\theta = \dfrac{\sec^6\theta}{32 - 48\sec^2\theta + 18\sec^4\theta - \sec^6\theta}\)} & A1~AG & {From the correct polynomial, reaches the printed result without error; a one-step multiplication is accepted.} \\
\nopagebreak & & [6] & \\
9(b) & {\(x = \sec\theta\):\\[4pt] \(\sec^6\theta - 18\sec^4\theta + 48\sec^2\theta - 32 = 0\)\\[4pt] \(\Rightarrow 32 - 48\sec^2\theta + 18\sec^4\theta - \sec^6\theta = 0\)} & M1 & {Substitutes \(x = \sec\theta\) and identifies the equation as the denominator in (a) equal to zero, sign change allowed.} \\
& {\(\sec 6\theta\) undefined \(\Rightarrow \cos 6\theta = 0\)} & A1 & {oe such as \(\sec 6\theta\) undefined or infinite.} \\
& {\(6\theta = \frac{\pi}{2}, \frac{3\pi}{2}, \ldots \Rightarrow \theta = \ldots\)} & M1 & {Solves their \(\cos 6\theta = 0\) for at least two values of \(\theta\) giving distinct values of \(\sec\theta\).} \\
& {\(x = \sec\frac{\pi}{12}, \sec\frac{\pi}{4}, \sec\frac{5\pi}{12}, \sec\frac{7\pi}{12}, \sec\frac{3\pi}{4}, \sec\frac{11\pi}{12}\)} & A1 & {Six distinct roots in the form \(\sec q\pi\), no others; any equivalent \(q\). Decimals, or \(\pm\sqrt2\), in place of \(\sec q\pi\) score A0.} \\
\nopagebreak & & [4] & \\
\nopagebreak \SetCell[c=2]{r} Total for Question 9 & & 10 & \\
\end{msscheme}

\msQuestion{10}
\begin{msscheme}{
  hline{3-7,9-13} = {2-Z}{0.4pt, msGrid},
  hline{8,14} = {0.4pt, msGrid},
  row{14} = {font=\bfseries},
}
Question & Scheme & Marks & Notes \\
10(a) & {\(\cos5\theta+i\sin5\theta=(\cos\theta+i\sin\theta)^5\)\\[4pt] \(\phantom{A}=\cos^5\theta+5i\cos^4\theta\sin\theta-10\cos^3\theta\sin^2\theta\)\\[4pt] \(\phantom{A=}-10i\cos^2\theta\sin^3\theta+5\cos\theta\sin^4\theta+i\sin^5\theta\)} & M1 & {Equates \(\cos5\theta+i\sin5\theta\) to \((\cos\theta+i\sin\theta)^5\) and attempts a binomial expansion with at least the three real terms. Condone coefficient or sign slips.} \\
& {\(\cos5\theta=\cos^5\theta-10\cos^3\theta\sin^2\theta+5\cos\theta\sin^4\theta\)} & A1 & {Correct real part; imaginary terms need not be seen.} \\
& {\(\cos5\theta=\cos^5\theta-10\cos^3\theta(1-\cos^2\theta)\)\\[4pt] \(\phantom{\cos5\theta=}+5\cos\theta(1-\cos^2\theta)^2\)} & M1 & {Substitutes \(\sin^2\theta=1-\cos^2\theta\) into their real part to give an expression in \(\cos\theta\) only.} \\
& {\(\phantom{A}=\cos^5\theta-10\cos^3\theta+10\cos^5\theta\)\\[4pt] \(\phantom{A=}+5\cos\theta-10\cos^3\theta+5\cos^5\theta\)} & dM1 & {Expands the brackets, including \((1-\cos^2\theta)^2\), and collects like powers of \(\cos\theta\). Condone slips.} \\
& {\(\cos5\theta=16\cos^5\theta-20\cos^3\theta+5\cos\theta\)} & A1~AG & {Correct real part leading without error to the printed result with \(\cos5\theta\) as LHS; combined steps accepted.} \\
\nopagebreak & & [5] & \\
10(b) & {\(\cos5\theta=\cos2\theta\)} & B1 & {Replaces the LHS by \(\cos5\theta\) using (a).} \\
& {\(5\theta=2n\pi\pm2\theta\Rightarrow\theta=\ldots\)} & M1 & {Solves by a general solution of \(\cos A=\cos B\) with both signs, a factor formula, or listing equal-cosine angles over a sufficient range, leading to \(\theta=\ldots\).} \\
& {\(\theta=\frac{2n\pi}{3}\)} & A1 & {Or its values \(0\) and \(\frac{2\pi}{3}\) in range.} \\
& {\(\theta=\frac{2n\pi}{7}\)} & A1 & {Or its values \(0\), \(\frac{2\pi}{7}\), \(\frac{4\pi}{7}\), \(\frac{6\pi}{7}\) in range.} \\
& {\(\theta=0,\ \frac{2\pi}{7},\ \frac{4\pi}{7},\ \frac{2\pi}{3},\ \frac{6\pi}{7}\)} & A1 & {All five exact values, including 0, and no others in \(0\le\theta\le\pi\).} \\
\nopagebreak & & [5] & \\
\nopagebreak \SetCell[c=2]{r} Total for Question 10 & & 10 & \\
\end{msscheme}

\msQuestion{11}
\begin{msscheme}{
  hline{3-7,9-12,14-16} = {2-Z}{0.4pt, msGrid},
  hline{8,13,17} = {0.4pt, msGrid},
  row{17} = {font=\bfseries},
}
Question & Scheme & Marks & Notes \\
11(a) & {\((\cos\theta+i\sin\theta)^6=\cos6\theta+i\sin6\theta\)\\[4pt] \(\phantom{(\cos\theta+i\sin\theta)^6}=\cos^6\theta+6i\cos^5\theta\sin\theta\)\\[4pt] \(\phantom{(\cos\theta+i\sin\theta)^6=}-15\cos^4\theta\sin^2\theta+\ldots\)} & M1 & {Equates \(\cos6\theta+i\sin6\theta\) to a binomial expansion of \((\cos\theta+i\sin\theta)^6\), at least its real terms, with powers of \(i\) applied. Condone coefficient and sign slips.} \\
& {\(\cos6\theta=\cos^6\theta-15\cos^4\theta\sin^2\theta\)\\[4pt] \(\phantom{\cos6\theta=}+15\cos^2\theta\sin^4\theta-\sin^6\theta\)} & A1 & \\
& {\(u=\sin^2\theta\): \(\cos6\theta=(1-u)^3-15(1-u)^2u\)\\[4pt] \(\phantom{u=\sin^2\theta: \cos6\theta=}+15(1-u)u^2-u^3\)} & M1 & {Replaces \(\cos^2\theta\) by \(1-\sin^2\theta\) throughout their real part.} \\
& {\(\phantom{\cos6\theta}=1-18u+48u^2-32u^3\)\\[4pt] \(\phantom{\cos6\theta}=(1-2u)(\ldots)\)} & dM1 & {Expands and collects into a cubic in \(\sin^2\theta\), then factorises it or expands the printed product to compare. Condone slips.} \\
& {\(\cos6\theta\equiv(1-2\sin^2\theta)(16\sin^4\theta\)\\[4pt] \(\phantom{\cos6\theta\equiv}-16\sin^2\theta+1)\)} & A1~AG & {Correct substituted expression leading without error to the printed identity; matching the expanded printed RHS with a conclusion is accepted.} \\
\nopagebreak & & [5] & \\
11(b) & {\((1-2u)(16u^2-16u+1)=0\Rightarrow u=\ldots\)} & M1 & {Sets each factor to zero and solves the quadratic in \(\sin^2\theta\) by formula with their values or completing the square.} \\
& {\(u=\frac12,\ \frac{2\pm\sqrt3}{4}\)} & A1 & {All three, oe such as \(\frac{16\pm8\sqrt3}{32}\).} \\
& {\(\cos\left(6\times\frac{\pi}{12}\right)=\cos\frac{\pi}{2}=0\)\\[4pt] \(0<\frac{\pi}{12}<\frac{\pi}{4}\Rightarrow\sin^2\frac{\pi}{12}<\frac12\)} & M1 & {Uses \(\theta=\frac{\pi}{12}\) as a solution of \(\cos6\theta=0\) and gives a reason rejecting \(\frac12\) and \(\frac{2+\sqrt3}{4}\), such as \(\sin^2\frac{\pi}{12}<\frac12\).} \\
& {\(\sin^2\left(\frac{\pi}{12}\right)=\frac{2-\sqrt3}{4}\)} & A1~AG & {cso, with both other values rejected.} \\
\nopagebreak & & [4] & \\
11(c) & {\(\sin^2\left(\frac{5\pi}{12}\right)=\frac{2+\sqrt3}{4}\)} & B1 & {Any exact form, here and below.} \\
& {\(\sin^2\left(\frac{7\pi}{12}\right)=\frac{2+\sqrt3}{4}\)} & B1 & \\
& {\(\sin^2\left(\frac{11\pi}{12}\right)=\frac{2-\sqrt3}{4}\)} & B1 & \\
\nopagebreak & & [3] & \\
\nopagebreak \SetCell[c=2]{r} Total for Question 11 & & 12 & \\
\end{msscheme}

\msQuestion{12}
\begin{msscheme}{
  hline{3-7,9-15} = {2-Z}{0.4pt, msGrid},
  hline{8,16} = {0.4pt, msGrid},
  row{16} = {font=\bfseries},
}
Question & Scheme & Marks & Notes \\
12(i) & {\(\cos 5\theta + i\sin 5\theta = (\cos\theta + i\sin\theta)^5\)\\[4pt] \(\phantom{\cos 5\theta} = \cos^5\theta + 5i\cos^4\theta\sin\theta\)\\[4pt] \(\phantom{\cos 5\theta =} - 10\cos^3\theta\sin^2\theta + \ldots\)} & M1 & {Uses de Moivre with a binomial expansion of \((\cos\theta + i\sin\theta)^5\), showing at least the real terms. Condone slips in powers of \(i\).} \\
& {\(\cos 5\theta = \cos^5\theta - 10\cos^3\theta\sin^2\theta\)\\[4pt] \(\phantom{\cos 5\theta =} + 5\cos\theta\sin^4\theta\)} & A1 & \\
& {\(\cos 5\theta = \cos\theta\big((1-\sin^2\theta)^2\)\\[4pt] \(\phantom{\cos 5\theta =} - 10(1-\sin^2\theta)\sin^2\theta + 5\sin^4\theta\big)\)} & M1 & {Replaces \(\cos^2\theta\) by \(1 - \sin^2\theta\) in their real part, keeping a factor \(\cos\theta\).} \\
& {\(\phantom{\cos 5\theta} = \cos\theta\big(1 - 2\sin^2\theta + \sin^4\theta\)\\[4pt] \(\phantom{\cos 5\theta =} - 10\sin^2\theta + 15\sin^4\theta\big)\)} & dM1 & {Expands and collects into powers of \(\sin\theta\). Condone arithmetic slips.} \\
& {\(\cos 5\theta \equiv \cos\theta(1 - 12\sin^2\theta + 16\sin^4\theta)\)} & A1~AG & {Correct substitution leading without error to the printed identity; expansion and collection may be combined.} \\
\nopagebreak & & [5] & \\
12(ii) & {\(\frac{\cos 5\theta}{\cos\theta} = 1 - 12u + 16u^2\), \(u = \sin^2\theta\)} & M1 & {Divides the identity by \(\cos\theta\) and treats the result as a quadratic in \(\sin^2\theta\), any letter.} \\
& {\(0 \le u < 1\)} & B1 & {Strict inequality at 1 needed; equivalent statements in \(\sin\theta\) accepted.} \\
& {\(u = \frac{12}{2 \times 16}\ [= \tfrac38]\)} & M1 & {Finds the turning point of their quadratic by differentiation, completing the square or \(-\frac{b}{2a}\). Condone sign slips.} \\
& {\(1 - 12\left(\tfrac38\right) + 16\left(\tfrac38\right)^2 = -\tfrac54\)} & A1 & {\(-\frac54\) oe.} \\
& {\(u \to 1 \Rightarrow 1 - 12 + 16 = \ldots\)} & M1 & {Evaluates their quadratic at, or as \(u\) tends to, \(u = 1\).} \\
& {\(5\), not attained: \(u = 1 \Rightarrow \cos\theta = 0\)} & A1 & {5 with the reason it is not attained. Solving \(f(u) = 5\) is not required.} \\
& {\(-\frac54 \le \frac{\cos 5\theta}{\cos\theta} < 5\)} & A1 & {Correct inequality at both ends; \(\left[-\frac54, 5\right)\) oe accepted.} \\
\nopagebreak & & [7] & \\
\nopagebreak \SetCell[c=2]{r} Total for Question 12 & & 12 & \\
\end{msscheme}

\msQuestion{13}
\begin{msscheme}{
  hline{3-7,9-12} = {2-Z}{0.4pt, msGrid},
  hline{8,13} = {0.4pt, msGrid},
  row{13} = {font=\bfseries},
}
Question & Scheme & Marks & Notes \\
13(a) & {\(\cos5\theta + i\sin5\theta = (c+is)^5\)\\[4pt] \(\phantom{\cos5\theta} = c^5 + 5ic^4s - 10c^3s^2 - 10ic^2s^3\)\\[4pt] \(\phantom{\cos5\theta =} + 5cs^4 + is^5\)} & M1 & {Equates \(\cos5\theta + i\sin5\theta\) to a binomial expansion of \((c+is)^5\) showing at least the real terms. Condone sign slips from powers of \(i\).} \\
& {\(\cos5\theta = c^5 - 10c^3s^2 + 5cs^4\)} & A1 & \\
& {\(\phantom{\cos5\theta} = c^5 - 10c^3(1-c^2) + 5c(1-c^2)^2\)} & M1 & {Replaces \(s^2\) by \(1-c^2\) throughout their real part, including \(s^4\) by \((1-c^2)^2\).} \\
& {\(\phantom{\cos5\theta} = c^5 - 10c^3 + 10c^5 + 5c - 10c^3 + 5c^5\)} & dM1 & {Multiplies out all brackets and collects like terms to a polynomial in \(\cos\theta\) only. Condone sign slips.} \\
& {\(\cos5\theta = 16\cos^5\theta - 20\cos^3\theta + 5\cos\theta\)} & A1~AG & {cso: correct substituted real part leading without error to the printed identity; combined steps accepted.} \\
\nopagebreak & & [5] & \\
13(b) & {\(x = \cos\theta\):\\[4pt] \(4(16x^5 - 20x^3 + 5x) + 1 = 4\cos5\theta + 1 = 0\)} & M1 & {Substitutes \(x = \cos\theta\) and uses part (a) to reach an equation in \(\cos5\theta\). Roots found without the identity score nothing in this part.} \\
& {\(\cos5\theta = -\tfrac14\)} & A1 & \\
& {\(5\theta = 104.48^\circ,\ 255.52^\circ\)\\[4pt] \(\Rightarrow \theta = 20.90^\circ,\ 51.10^\circ \Rightarrow x = \cos\theta = \ldots\)} & M1 & {Finds \(\arccos\) of their \(\cos5\theta\) and a second solution such as \(360^\circ - \alpha\) (or \(2\pi - \alpha\)), divides by 5 and takes cosines.} \\
& {\(x = 0.934,\ 0.628\)} & A1 & {awrt 0.934 and 0.628, with no other positive roots given.} \\
\nopagebreak & & [4] & \\
\nopagebreak \SetCell[c=2]{r} Total for Question 13 & & 9 & \\
\end{msscheme}

\msQuestion{14}
\begin{msscheme}{
  hline{3-7,9-14} = {2-Z}{0.4pt, msGrid},
  hline{8,15} = {0.4pt, msGrid},
  row{15} = {font=\bfseries},
}
Question & Scheme & Marks & Notes \\
14(a) & {\(\cos4\theta + i\sin4\theta = (\cos\theta+i\sin\theta)^4\)\\[4pt] \(\phantom{A} = \cos^4\theta + 4i\cos^3\theta\sin\theta - 6\cos^2\theta\sin^2\theta\)\\[4pt] \(\phantom{A =} - 4i\cos\theta\sin^3\theta + \sin^4\theta\)} & M1 & {Equates \(\cos4\theta + i\sin4\theta\) to a binomial expansion of \((\cos\theta+i\sin\theta)^4\) showing at least the three real terms. Condone slips in coefficients or powers of \(i\).} \\
& {\(\cos4\theta = \cos^4\theta - 6\sin^2\theta\cos^2\theta + \sin^4\theta\)} & A1 & \\
& {\(\cos4\theta = (\sin^2\theta+\cos^2\theta)^2 - 8\sin^2\theta\cos^2\theta\)} & M1 & {Uses \(\sin^2\theta+\cos^2\theta=1\) to rewrite the fourth powers, e.g. completing the square or substituting \(\cos^2\theta = 1-\sin^2\theta\).} \\
& {\(\cos4\theta = 1 - 8\sin^2\theta\cos^2\theta\)} & dM1 & {Collects to the form \(\cos4\theta = 1 - k\sin^2\theta\cos^2\theta\).} \\
& {\(\sin^2\theta\cos^2\theta = \frac18(1-\cos4\theta)\)} & A1~AG & {Printed result reached without error from a correct real-part equation; combined steps accepted. Nothing for working back from the result.} \\
\nopagebreak & & [5] & \\
14(b) & {\(e^{\int -2\tan\theta\,\mathrm{d}\theta} = e^{-2\ln\sec\theta} = \ldots\)} & M1 & {Forms \(e^{\int \pm 2\tan\theta\,\mathrm{d}\theta}\) and integrates to the form \(e^{k\ln\sec\theta}\) or \(e^{k\ln\cos\theta}\).} \\
& {\(\phantom{A} = \cos^2\theta\)} & A1 & \\
& {\(\frac{\mathrm{d}}{\mathrm{d}\theta}(y\cos^2\theta) = \sin^2\theta\cos^2\theta = \frac18(1-\cos4\theta)\)\\[4pt] \(\Rightarrow y\cos^2\theta = \ldots\)} & dM1 & {Multiplies through by their integrating factor, writes LHS as \(\frac{\mathrm{d}}{\mathrm{d}\theta}(y \times\) their factor\()\) and integrates the RHS, using (a) or otherwise, with at least one term correct.} \\
& {\(y\cos^2\theta = \frac{\theta}{8} - \frac{\sin4\theta}{32} + c\)} & A1 & {oe; needs the constant.} \\
& {\(0 = \frac{\pi}{32} - \frac{\sin\pi}{32} + c \Rightarrow c = \ldots\ [= -\frac{\pi}{32}]\)} & M1 & {Substitutes \(\theta = \frac{\pi}{4}\), \(y = 0\) into their general solution containing a constant and solves for it.} \\
& {\(y = \frac{4\theta - \sin4\theta - \pi}{32\cos^2\theta}\)} & A1 & {oe; \(y\cos^2\theta = \frac{\theta}{8} - \frac{\sin4\theta}{32} - \frac{\pi}{32}\) accepted.} \\
\nopagebreak & & [6] & \\
\nopagebreak \SetCell[c=2]{r} Total for Question 14 & & 11 & \\
\end{msscheme}

\msQuestion{15}
\begin{msscheme}{
  hline{3-6,8-9} = {2-Z}{0.4pt, msGrid},
  hline{7,10} = {0.4pt, msGrid},
  row{10} = {font=\bfseries},
}
Question & Scheme & Marks & Notes \\
15(a) & {\((z+z^{-1})^6 = 64\cos^6\theta\)} & B1 & \\
& {\((z+z^{-1})^6 = z^6+6z^4+15z^2+20\)\\[4pt] \(\phantom{(z+z^{-1})^6 =} +15z^{-2}+6z^{-4}+z^{-6}\)} & M1 A1 & {\msNote{M1}{Attempts the binomial expansion: seven terms with binomial coefficients and powers of \(z\). Powers need not be simplified; condone one coefficient slip.}\\ \msNote{A1}{Correct expansion with powers simplified.}} \\
& {\(64\cos^6\theta = (z^6+z^{-6})+6(z^4+z^{-4})+\ldots\)\\[4pt] \(\phantom{64\cos^6\theta} = 2\cos6\theta+12\cos4\theta+30\cos2\theta+20\)} & M1 & {Pairs \(z^n\) with \(z^{-n}\) and replaces by \(2\cos n\theta\) for at least two pairs, with their coefficients.} \\
& {\(a=\frac{1}{32},\ b=\frac{3}{16},\ c=\frac{15}{32},\ d=\frac{5}{16}\)} & A1 & {All four, oe, stated or in the identity.} \\
\nopagebreak & & [5] & \\
15(b) & {\(\frac{1}{32}\int_0^{\pi/4}(\cos6\theta+6\cos4\theta+15\cos2\theta+10)\,\mathrm{d}\theta\)\\[4pt] \(= \frac{1}{32}\left[\frac16\sin6\theta+\frac32\sin4\theta+\frac{15}{2}\sin2\theta+10\theta\right]_0^{\pi/4}\)} & M1 A1 & {\msNote{M1}{Integrates their part (a) identity term by term to the form \(p\sin6\theta+q\sin4\theta+r\sin2\theta+s\theta\). Condone sign and coefficient slips.}\\ \msNote{A1}{Correct integral, oe.}} \\
& {\(= \frac{1}{32}\left(-\frac16+0+\frac{15}{2}+\frac{5\pi}{2}\right)-0\)\\[4pt] \(= \frac{11}{48}+\frac{5\pi}{64}\)} & dM1 A1~AG & {\msNote{dM1}{Substitutes \(\theta=\frac{\pi}{4}\) and \(\theta=0\) and subtracts; the zero lower limit may be implied.}\\ \msNote{A1~AG}{Correct substituted expression with sine values evaluated, leading without error to the printed result; combined steps accepted.}} \\
\nopagebreak & & [4] & \\
\nopagebreak \SetCell[c=2]{r} Total for Question 15 & & 9 & \\
\end{msscheme}

\msQuestion{16}
\begin{msscheme}{
  hline{3-4,6-7,9-13,15-18,20-22} = {2-Z}{0.4pt, msGrid},
  hline{5,8,14,19,23} = {0.4pt, msGrid},
  row{23} = {font=\bfseries},
}
Question & Scheme & Marks & Notes \\
16(i)(A) & {\((\cos\theta+i\sin\theta)^n = (e^{i\theta})^n = e^{in\theta}\)} & M1 & {Raises both sides of Euler's formula to the power \(n\) and writes \((e^{i\theta})^n = e^{in\theta}\).} \\
& {\(e^{in\theta} = \cos n\theta + i\sin n\theta\)} & A1~AG & {Applies Euler's formula with \(n\theta\) to reach the printed result. Quoting de Moivre's theorem earns nothing.} \\
\nopagebreak & & [2] & \\
16(i)(B) & {\(e^{-i\theta} = \cos(-\theta)+i\sin(-\theta)\)\\[4pt] \(\phantom{e^{-i\theta}} = \cos\theta - i\sin\theta\)} & M1 & {Obtains \(e^{-i\theta}\) in terms of \(\cos\theta\) and \(\sin\theta\), e.g. by replacing \(\theta\) by \(-\theta\), and subtracts the two results. Condone sign slips.} \\
& {\(e^{i\theta}-e^{-i\theta} = 2i\sin\theta\)\\[4pt] \(\Rightarrow \sin\theta = \frac{1}{2i}(e^{i\theta}-e^{-i\theta})\)} & A1~AG & {Correct \(e^{-i\theta} = \cos\theta - i\sin\theta\), with \(\cos(-\theta)=\cos\theta\) and \(\sin(-\theta)=-\sin\theta\) used, leading without error to the printed result.} \\
\nopagebreak & & [2] & \\
16(ii) & {\(\cos 8\theta + i\sin 8\theta = (\cos\theta+i\sin\theta)^8\)\\[4pt] \(\phantom{\cos 8\theta} = \cos^8\theta + 8i\cos^7\theta\sin\theta\)\\[4pt] \(\phantom{\cos 8\theta =} - 28\cos^6\theta\sin^2\theta + \ldots\)} & M1 & {Uses de Moivre with \(n=8\) and expands \((\cos\theta+i\sin\theta)^8\) binomially, at least the real terms. Condone coefficient or sign slips.} \\
& {\(\cos 8\theta = \cos^8\theta - 28\cos^6\theta\sin^2\theta\)\\[4pt] \(\phantom{\cos 8\theta =} + 70\cos^4\theta\sin^4\theta - 28\cos^2\theta\sin^6\theta\)\\[4pt] \(\phantom{\cos 8\theta =} + \sin^8\theta\)} & A1 & \\
& {\(\cos 8\theta = (1-\sin^2\theta)^4\)\\[4pt] \(\phantom{\cos 8\theta =} - 28(1-\sin^2\theta)^3\sin^2\theta + \ldots\)} & M1 & {Replaces every \(\cos^2\theta\) in their real part by \(1-\sin^2\theta\).} \\
& {\(\cos 8\theta = (1 - 4\sin^2\theta + 6\sin^4\theta - \ldots)\)\\[4pt] \(\phantom{\cos 8\theta =} - 28(\sin^2\theta - 3\sin^4\theta + \ldots) + \ldots\)} & dM1 & {Expands all the brackets and collects like powers of \(\sin\theta\). Condone sign and arithmetic slips.} \\
& {\(a=128,\ b=-256,\ c=160,\ d=-32,\ e=1\)} & A1 A1 & {\msNote{A1}{At least three of the five constants correct.}\\ \msNote{A1}{All five correct; the identity with these coefficients is accepted.}} \\
\nopagebreak & & [6] & \\
16(iii) & {\(\sin^8\theta = \frac{1}{(2i)^8}(e^{i\theta}-e^{-i\theta})^8\)\\[4pt] \(\phantom{\sin^8\theta} = \frac{1}{256}(e^{i\theta}-e^{-i\theta})^8\)} & B1 & {Needs \((2i)^8\) evaluated as 256.} \\
& {\(\phantom{\sin^8\theta} = \frac{1}{256}(e^{8i\theta} - 8e^{6i\theta} + 28e^{4i\theta} - 56e^{2i\theta} + 70\)\\[4pt] \(\phantom{\sin^8\theta =} - 56e^{-2i\theta} + 28e^{-4i\theta} - 8e^{-6i\theta} + e^{-8i\theta})\)} & M1 A1 & {\msNote{M1}{Expands the eighth power binomially with coefficients from row 8 of Pascal's triangle, combining powers of \(e^{i\theta}\). Accept \(z = e^{i\theta}\) notation. Condone sign slips.}\\ \msNote{A1}{All nine terms correct.}} \\
& {\(\sin^8\theta = \frac{1}{256}(2\cos 8\theta - 16\cos 6\theta + 56\cos 4\theta\)\\[4pt] \(\phantom{\sin^8\theta =} - 112\cos 2\theta + 70)\)} & M1 & {Pairs \(e^{ik\theta}\) with \(e^{-ik\theta}\) as \(2\cos k\theta\) for at least two values of \(k\).} \\
& {\(P=\frac{1}{128},\ Q=-\frac{1}{16},\ R=\frac{7}{32},\)\\[4pt] \(S=-\frac{7}{16},\ T=\frac{35}{128}\)} & A1 A1 & {\msNote{A1}{At least three of the five constants correct.}\\ \msNote{A1}{All five correct, oe.}} \\
\nopagebreak & & [6] & \\
16(iv) & {\(\sin^8\frac{\pi}{8} = \frac{1}{128}\left(-1 - 8\left(-\frac{\sqrt2}{2}\right) + 28(0)\right.\)\\[4pt] \(\phantom{\sin^8\frac{\pi}{8} =} \left. - 56\left(\frac{\sqrt2}{2}\right) + 35\right)\)} & M1 & {Substitutes \(\theta=\frac{\pi}{8}\) into their identity of the form in (iii) and uses exact cosine values. Decimal cosines do not score.} \\
& {\(\sin^8\frac{\pi}{8} = \frac{34-24\sqrt2}{128} = \frac{17-12\sqrt2}{64}\)} & A1 & {cso, from correct constants in (iii).} \\
& {\(0<\frac{\pi}{8}<\frac{\pi}{2} \Rightarrow \sin\frac{\pi}{8}>0\)\\[4pt] \(\Rightarrow \sin\frac{\pi}{8} = \left(\frac{17-12\sqrt2}{64}\right)^{\frac18}\)} & A1~AG & {Needs a reason for taking the positive real eighth root, such as \(\sin\frac{\pi}{8}>0\) or \(\frac{\pi}{8}\) acute.} \\
\nopagebreak & & [3] & \\
\nopagebreak \SetCell[c=2]{r} Total for Question 16 & & 19 & \\
\end{msscheme}

\msQuestion{17}
\begin{msscheme}{
  hline{3-13,15-17,19-23} = {2-Z}{0.4pt, msGrid},
  hline{14,18,24} = {0.4pt, msGrid},
  cell{7-12}{2-4} = {bg=msTint},
  row{24} = {font=\bfseries},
}
Question & Scheme & Marks & Notes \\
17(a) & {\((z-z^{-1})^4(z+z^{-1})^2 = (2i\sin\theta)^4(2\cos\theta)^2\)\\[4pt] \(\phantom{(z-z^{-1})^4(z+z^{-1})^2} = 64\sin^4\theta\cos^2\theta\)} & B1 & \\
& {\((z-z^{-1})^4(z+z^{-1})^2 = (z^2-2+z^{-2})(z^4-2+z^{-4})\)\\[4pt] \(\phantom{(z-z^{-1})^4(z+z^{-1})^2} = \ldots\)} & M1 & {Attempts the full expansion of the product into powers of \(z\), in any order. Condone sign and coefficient slips.} \\
& {\(= z^6-2z^4-z^2+4-z^{-2}-2z^{-4}+z^{-6}\)} & A1 & {Correct seven-term expansion, oe.} \\
& {\(= (z^6+z^{-6})-2(z^4+z^{-4})-(z^2+z^{-2})+4\)\\[4pt] \(= 2\cos 6\theta-4\cos 4\theta-2\cos 2\theta+4\)} & M1 & {Pairs \(z^n\) with \(z^{-n}\) and replaces each pair by \(2\cos n\theta\).} \\
& {\(64\sin^4\theta\cos^2\theta = 2\cos 6\theta-4\cos 4\theta-2\cos 2\theta+4\)\\[4pt] \(\Rightarrow \sin^4\theta\cos^2\theta \equiv \frac{1}{32}(\cos 6\theta-2\cos 4\theta-\cos 2\theta+2)\)} & A1~AG & {Correct expansion equated to \(64\sin^4\theta\cos^2\theta\), leading without error to the printed identity.} \\
& \SetCell[c=3]{l} \msAlternative{Double-angle formulae in \(\cos 2\theta\)} & & \\
\nopagebreak & {\(\sin^4\theta\cos^2\theta = \left(\frac{1-\cos 2\theta}{2}\right)^2\left(\frac{1+\cos 2\theta}{2}\right) = \ldots\)} & M1 & {Uses \(\sin^2\theta=\frac12(1-\cos 2\theta)\), \(\cos^2\theta=\frac12(1+\cos 2\theta)\) and expands to a polynomial in \(\cos 2\theta\). Condone sign slips.} \\
& {\(= \frac18(1-\cos 2\theta-\cos^2 2\theta+\cos^3 2\theta)\)} & A1 & \\
& {\((\cos 2\theta+i\sin 2\theta)^3 = \cos 6\theta+i\sin 6\theta\)\\[4pt] \(\Rightarrow \cos 6\theta = \cos^3 2\theta-3\cos 2\theta\sin^2 2\theta = \ldots\)} & M1 & {Expands \((\cos 2\theta+i\sin 2\theta)^3\) and equates real parts to express \(\cos 6\theta\) in powers of \(\cos 2\theta\); \(\cos 4\theta\) may come from a double-angle formula. Not for quoting the triple-angle formula.} \\
& {\(\cos^2 2\theta = \frac12(1+\cos 4\theta)\)\\[4pt] \(\cos^3 2\theta = \frac14(\cos 6\theta+3\cos 2\theta)\)} & A1 & {Both, oe.} \\
& {\(\sin^4\theta\cos^2\theta = \frac18\left(1-\cos 2\theta-\frac{1+\cos 4\theta}{2}+\frac{\cos 6\theta+3\cos 2\theta}{4}\right)\)\\[4pt] \(\phantom{\sin^4\theta\cos^2\theta} \equiv \frac{1}{32}(\cos 6\theta-2\cos 4\theta-\cos 2\theta+2)\)} & A1~AG & {Correct substitution into the correct cubic, leading without error to the printed identity.} \\
\nopagebreak & & [5] & \\
17(b) & {\(\sin^2\theta\cos^4\theta = \sin^4\alpha\cos^2\alpha\)\\[4pt] \(\phantom{\sin^2\theta\cos^4\theta} = \frac{1}{32}(\cos 6\alpha-2\cos 4\alpha-\cos 2\alpha+2)\)} & M1 & {Replaces \(\sin\theta\) by \(\cos\alpha\) and \(\cos\theta\) by \(\sin\alpha\) and applies the identity of (a) in \(\alpha\).} \\
& {\(\cos 2\alpha = -\cos 2\theta,\ \cos 4\alpha = \cos 4\theta,\ \cos 6\alpha = -\cos 6\theta\)} & M1 & {Uses \(n\alpha=\frac{n\pi}{2}-n\theta\) to write all three cosines in terms of \(\theta\). Condone sign slips.} \\
& {\(\sin^2\theta\cos^4\theta \equiv \frac{1}{32}(-\cos 6\theta-2\cos 4\theta+\cos 2\theta+2)\)} & A1 & {oe.} \\
\nopagebreak & & [3] & \\
17(c) & {\(\frac{1}{32}(\cos 6\theta-2\cos 4\theta-\cos 2\theta+2)\)\\[4pt] \(\phantom{=} + \frac{2}{32}(\their{-\cos 6\theta-2\cos 4\theta+\cos 2\theta+2})\)} & M1 & {Adds the identity from (a) to twice their identity from (b), of the form \(\frac{1}{32}(\)cosines of \(2\theta, 4\theta, 6\theta\) and a constant\()\).} \\
& {\(= \frac{1}{32}(-\cos 6\theta-6\cos 4\theta+\cos 2\theta+6)\)} & A1 & \\
& {\(\frac{1}{32}\left[-\frac16\sin 6\theta-\frac32\sin 4\theta+\frac12\sin 2\theta+6\theta\right]_0^a\)} & M1 & {Integrates at least one cosine term to the form \(k\sin n\theta\) and substitutes the limits 0 and \(a\).} \\
& {\(\frac{1}{32}\left(6a+\frac12\sin 2a-\frac32\sin 4a-\frac16\sin 6a\right)\)\\[4pt] \(\phantom{=} = \frac{\pi}{32}-\frac{\sqrt3}{64}\)} & A1 & {Correct expression in \(a\), oe, equated to the given value.} \\
& {\(a = \frac{\pi}{6}\)} & A1 & {Exact; may be found by trial and checked by substitution. Uniqueness need not be shown.} \\
\nopagebreak & & [5] & \\
\nopagebreak \SetCell[c=2]{r} Total for Question 17 & & 13 & \\
\end{msscheme}

\msQuestion{18}
\begin{msscheme}{
  hline{3-11,13-17,19-21} = {2-Z}{0.4pt, msGrid},
  hline{12,18,22} = {0.4pt, msGrid},
  cell{6-10}{2-4} = {bg=msTint},
  row{22} = {font=\bfseries},
}
Question & Scheme & Marks & Notes \\
18(a) & {\(z = \cos\theta + i\sin\theta \Rightarrow (z + z^{-1})^5 = 32\cos^5\theta\)} & B1 & \\
& {\((z + z^{-1})^5 = z^5 + 5z^3 + 10z + 10z^{-1}\)\\[4pt] \(\phantom{(z + z^{-1})^5 =} + 5z^{-3} + z^{-5}\)} & M1 A1 & {\msNote{M1}{Expands to six terms in \(z^5, z^3, \ldots, z^{-5}\) with binomial coefficients. Condone coefficient slips.}\\ \msNote{A1}{Correct expansion.}} \\
& {\(\phantom{(z + z^{-1})^5} = 2\cos 5\theta + 10\cos 3\theta + 20\cos\theta\)} & M1 & {Pairs terms as \(z^n + z^{-n}\) and replaces each pair by \(2\cos n\theta\).} \\
& {\(A = \frac1{16},\ B = \frac5{16},\ C = \frac58\)} & A1 & {Values stated, or the identity written with them.} \\
& \SetCell[c=3]{l} \msAlternative{Multiple-angle polynomials} & & \\
\nopagebreak & {\(\cos 5\theta = \mathrm{Re}(\cos\theta + i\sin\theta)^5\)\\[4pt] \(\phantom{\cos 5\theta} = \cos^5\theta - 10\cos^3\theta\sin^2\theta + 5\cos\theta\sin^4\theta\)\\[4pt] \(\cos 3\theta = \cos^3\theta - 3\cos\theta\sin^2\theta\)} & M1 & {Expands \((\cos\theta + i\sin\theta)^n\) and equates real parts, for \(n = 5\) at least; \(\cos 3\theta = 4\cos^3\theta - 3\cos\theta\) may be quoted. Condone coefficient slips.} \\
& {\(\cos 5\theta = 16\cos^5\theta - 20\cos^3\theta + 5\cos\theta\)\\[4pt] \(\cos 3\theta = 4\cos^3\theta - 3\cos\theta\)} & A1 & {Both correct, using \(\sin^2\theta = 1 - \cos^2\theta\).} \\
& {\(16A = 1,\ -20A + 4B = 0,\ 5A - 3B + C = 0\)} & M1 & {Substitutes their polynomials into the identity and equates coefficients of at least two powers of \(\cos\theta\).} \\
& {\(A = \frac1{16},\ B = \frac5{16},\ C = \frac58\)} & dM1 A1 & {\msNote{dM1}{Solves their equations for \(A\), \(B\) and \(C\).}\\ \msNote{A1}{All three correct.}} \\
\nopagebreak & & [5] & \\
18(b) & {\(\theta = \cos^{-1}(x^{1/6})\)\\[4pt] \(f'(x) = -(15\cos 5\theta + 75\cos 3\theta + 150\cos\theta)\frac{\mathrm{d}\theta}{\mathrm{d}x}\)} & M1 A1 & {\msNote{M1}{Differentiates each \(\sin k\theta\) to \(\pm k\cos k\theta\) times \(\frac{\mathrm{d}\theta}{\mathrm{d}x}\) or an inner derivative. Condone coefficient slips.}\\ \msNote{A1}{Correct outer derivatives; the inner derivative may be unevaluated or incorrect.}} \\
& {\(\frac{\mathrm{d}\theta}{\mathrm{d}x} = -\frac{1}{\sqrt{1 - (x^{1/6})^2}} \times \frac16 x^{-5/6}\)} & M1 & {Uses \(\frac{\mathrm{d}}{\mathrm{d}u}\cos^{-1}u = \pm\frac{1}{\sqrt{1-u^2}}\) with \(\frac{\mathrm{d}u}{\mathrm{d}x} = kx^{-5/6}\), or differentiates \(\cos\theta = x^{1/6}\) implicitly.} \\
& {\(\phantom{\frac{\mathrm{d}\theta}{\mathrm{d}x}} = -\frac{x^{-5/6}}{6\sqrt{1 - x^{1/3}}}\)} & A1 & {oe, in terms of \(x\).} \\
& {\(15(\cos 5\theta + 5\cos 3\theta + 10\cos\theta) = 15 \times 16\cos^5\theta\)} & M1 & {Uses part (a) to replace the cosine sum by a multiple of \(\cos^5\theta\).} \\
& {\(f'(x) = 240x^{5/6} \times \frac{x^{-5/6}}{6\sqrt{1 - x^{1/3}}}\)\\[4pt] \(\phantom{f'(x)} = \frac{40}{\sqrt{1 - x^{1/3}}}\)} & A1~AG & {Correct substituted expression with \(\cos^5\theta = x^{5/6}\), leading without error to the printed result.} \\
\nopagebreak & & [6] & \\
18(c) & {Area \(= \frac1{40}\int_0^1 f'(x)\,\mathrm{d}x = \frac1{40}\big[f(x)\big]_0^1\)} & M1 & {Uses part (b) to write the area as \(\frac1{40}\) of \(f\) between 0 and 1. Limit notation at 1 need not be shown.} \\
& {\(x \to 1 \Rightarrow f(x) \to 0\); \(f(0) = -3 + 25 - 150 = -128\)\\[4pt] \(\Rightarrow \frac1{40}\big(0 - (-128)\big)\)} & dM1 & {Uses \(\cos^{-1}1 = 0\) and \(\cos^{-1}0 = \frac{\pi}{2}\) to evaluate \(f\) at both limits and subtracts in the correct order. Condone one sign slip in \(f(0)\).} \\
& {\(\frac{16}{5}\)} & A1 & {cso; 3.2 accepted.} \\
\nopagebreak & & [3] & \\
\nopagebreak \SetCell[c=2]{r} Total for Question 18 & & 14 & \\
\end{msscheme}

\end{document}
